\chapter{Conclusão}\label{cap:conclusao}

\section{Principais Contribuições}
Como orientado no Manual do SiBI, nessa seção o pesquisador retoma seus objetivos de pesquisa e suas hipóteses / suposições a fim de apresentar os principais resultados obtidos a partir da análise dos dados. Na seção anterior, o investigador fez uma análise minuciosa; aqui, resume os principais pontos, associando-os com os objetivos que pretendia alcançar. 

\section{Limitações da Pesquisa}
O pesquisador deve se antecipar e relatar os problemas encontrados na condução de sua pesquisa. Esse relato, além de mostrar o rigor metodológico e o compromisso com a busca da verdade, permite que o leitor entenda com mais clareza as decisões que nortearam a investigação.

\section{Pesquisas Futuras}\index{pesquisas futuras}
O pesquisador pode usar como inspiração para novas perguntas de pesquisa as limitações a que sua pesquisa esteve sujeita, ou seja, focar no que poderia ter sido feito de forma diferente para melhorar a coleta ou a análise dos dados. Ou qual outro método de pesquisa poderia ter sido utilizado e com quais objetivos.
Pode ainda, a partir dos resultados obtidos em sua investigação, propor novas questões acerca do fenômeno em estudo, estimulando a continuidade das pesquisas sobre o tema.